\documentclass[12pt]{article}
\usepackage[utf8]{inputenc}
\usepackage[T1,T2A]{fontenc}
\usepackage[english,russian]{babel}
\usepackage{graphicx}
\usepackage{subcaption}
\usepackage{amsmath}
\usepackage{fancyhdr}
\usepackage{geometry}
\usepackage{dirtytalk}
\usepackage{csquotes}
\usepackage{hyperref}
\usepackage{listings}
\usepackage{xcolor}
\usepackage{array}
\usepackage{caption}
\usepackage{float}
\usepackage{booktabs}
\usepackage{tikz}
\usepackage{pgfgantt}
\usetikzlibrary{shapes,arrows,positioning,fit,backgrounds,automata,calc}

% Code styling
\lstset{
    basicstyle=\ttfamily\footnotesize,
    backgroundcolor=\color{lightgray!20},
    keywordstyle=\color{blue},
    commentstyle=\color{green!50!black},
    stringstyle=\color{red},
    numbers=left,
    numberstyle=\tiny\color{gray},
    stepnumber=1,
    numbersep=10pt,
    showspaces=false,
    showstringspaces=false,
    showtabs=false,
    frame=single,
    tabsize=2,
    breaklines=true,
    breakatwhitespace=false,
    escapeinside={\%*}{*)},
    extendedchars=false,
    keepspaces=true,
    morekeywords={void, TaskHandle_t, TickType_t, xTaskCreate, vTaskDelay,
                  vTaskDelayUntil, xSemaphoreGive, xSemaphoreTake,
                  xSemaphoreCreateMutex, xSemaphoreCreateBinary, xTaskGetTickCount,
                  pdMS_TO_TICKS, tskIDLE_PRIORITY, UBaseType_t}
}

\linespread{1.25}
\setlength{\parindent}{0.8cm}
\setlength{\parskip}{0em}
\renewcommand{\headrulewidth}{0pt}
\geometry{a4paper, portrait, margin=1in}
\setlength{\headheight}{14.49998pt}

\begin{document}
\begin{titlepage}
    \thispagestyle{empty}

    \begin{center}
        \textsc{\large Ministry of Education of Republic of Moldova}\\[0.5cm]
        \textsc{\large Technical University of Moldova}\\[0.5cm]
        \textsc{\large Faculty of Computers, Informatics and Microelectronics}\\[0.5cm]
        \textsc{\large Department of Physics}\\[1.2cm]
    \end{center}

    \vspace{25mm}

    \begin{center}
        \textsc{\Large Embedded Systems}\\[0.5cm]
        \textsc{\large Laboratory work \#5.1}\\[0.5cm]

        \newcommand{\HRule}{\rule{\linewidth}{0.5mm}}
        \vspace{10mm}
        \HRule \\[0.4cm]
        { \LARGE \bfseries ON-OFF Temperature Control System with Hysteresis}\\[0.4cm]
        \HRule \\[1.5cm]
    \end{center}

    \vspace{10mm}

    \begin{minipage}[t]{0.4\textwidth}
        \begin{flushleft} \large
            \emph{Author:} \\
            Dmitrii \textsc{Belih}\\
            std. gr. FAF-232
        \end{flushleft}
    \end{minipage}
    \hfill
    \begin{minipage}[t]{0.4\textwidth}
        \begin{flushright} \large
            \emph{Verified:} \\
            \textsc{Martiniuc} A.\\
        \end{flushright}
    \end{minipage}\\[3cm]

    \vspace{5mm}
    \begin{center}
        \large Chisinau 2026 \\[0.5cm]
    \end{center}

    \vfill
\end{titlepage}

\setcounter{page}{1}
\pagestyle{fancy}
\fancyhf{}
\rhead{\thepage}
\lhead{FAF-232 Belih Dmitrii; Laboratory Work 5.1}

% ============================================================================
% CHAPTER 1: INTRODUCTION
% ============================================================================
\section*{1. Introduction}

\subsection*{1.1. Purpose of the Laboratory Work}

The purpose of this laboratory work is to design and implement a modular, real-time
embedded application for an ESP32 microcontroller that performs ON-OFF temperature
control with hysteresis (Variant~A). The system reads temperature from a DHT11
digital sensor, compares it with a configurable Set Point, and actuates a relay to
maintain the temperature within the desired range. Hysteresis is applied to the
control law to prevent rapid relay toggling caused by small fluctuations around the
Set Point. The implementation is structured as a multi-task FreeRTOS application in
which acquisition, control, and display run as independent concurrent tasks with
justified recurrence periods. The user can adjust the Set Point interactively through
a push button (short press: $+1\,^{\circ}$C; long press: $-1\,^{\circ}$C) or via
serial commands. Measured temperature, Set Point, and relay state are displayed on a
16$\times$2 I2C LCD and streamed to the Arduino Serial Plotter for real-time
visualization.

\subsection*{1.2. Laboratory Objectives}

The specific objectives of this laboratory work are:
\begin{itemize}
    \item Understand the ON-OFF control strategy with two-point hysteresis and its
          impact on actuator switching frequency.
    \item Implement a FreeRTOS multi-task architecture with acquisition, control, and
          display tasks running at independent recurrence periods.
    \item Integrate a DHT11 digital temperature sensor using the Adafruit DHT library.
    \item Implement thread-safe access to shared peripherals (LCD) using a FreeRTOS
          mutex and coordinate task signalling via a binary semaphore.
    \item Provide interactive Set Point configuration through a button interface and a
          serial command parser.
    \item Stream real-time data (SetPoint, Temp, Output) to the Arduino Serial Plotter.
    \item Validate the hysteresis deadband behaviour and verify that no relay toggling
          occurs within the configured deadband.
\end{itemize}

% ============================================================================
% CHAPTER 2: DOMAIN ANALYSIS
% ============================================================================
\section*{2. Analysis of Domain Situation and Technological Context}

\subsection*{2.1. ON-OFF Control and Hysteresis}

ON-OFF (bang-bang) control is the simplest closed-loop strategy: the actuator is
either fully ON or fully OFF depending on whether the process variable is below or
above the Set Point. Without hysteresis, any noise on the measured value around the
Set Point causes extremely frequent switching (chattering), which shortens actuator
life. Hysteresis introduces a deadband $[SP - H,\; SP + H]$ where the current relay
state is preserved:

\begin{equation}
    \text{Relay} =
    \begin{cases}
        \text{ON}                     & \text{if } T < SP - H \quad\text{(below lower bound)}\\
        \text{OFF}                    & \text{if } T > SP + H \quad\text{(above upper bound)}\\
        \text{unchanged (hold state)} & \text{if } SP - H \le T \le SP + H
    \end{cases}
\end{equation}

The total deadband width is $2H$. A larger $H$ means fewer switch events but a wider
temperature oscillation band; a smaller $H$ tightens regulation but increases switching
frequency.

\subsection*{2.2. DHT11 Digital Temperature and Humidity Sensor}

The DHT11 is a low-cost digital sensor that uses a single-wire serial protocol to
transmit a 40-bit frame containing 8-bit integer and fractional parts for both
humidity and temperature. The minimum sampling interval is 1\,s, which sets the lower
bound on the acquisition task period. The Adafruit DHT library abstracts the protocol
and exposes a simple \texttt{readTemperature()} API. If no sensor is detected
(returned value $<-100\,^{\circ}$C), the system falls back to a software simulation:
temperature rises by $0.2\,^{\circ}$C per acquisition cycle when the relay is ON, and
falls by $0.1\,^{\circ}$C when the relay is OFF, enabling demonstration without
physical hardware.

\subsection*{2.3. Relay Actuation}

A relay module provides galvanic isolation between the low-voltage ESP32 control
signal and the load circuit. The relay coil is driven from the 5\,V rail through an
opto-coupler and a transistor driver built into the module. Driving GPIO~23 HIGH
energises the coil and closes the normally-open (NO) contact; driving it LOW
de-energises the coil and opens the contact. The binary \texttt{Actuator} driver
wraps \texttt{digitalWrite} calls and tracks the relay state and last-toggle
timestamp for wear monitoring.

\subsection*{2.4. FreeRTOS for Embedded Control Systems}

FreeRTOS provides deterministic preemptive scheduling that is essential for
closed-loop control. The use of \texttt{vTaskDelayUntil()} instead of
\texttt{vTaskDelay()} ensures drift-free periodic execution: the wake-up time is
advanced by an exact tick count on every iteration, so accumulated jitter is
bounded. Priority-based scheduling ensures that the higher-priority control tasks
(P3) always pre-empt the lower-priority display task (P2), guaranteeing control
loop responsiveness.

\subsection*{2.5. Inter-Task Communication and Synchronisation}

Two synchronisation primitives are used:
\begin{itemize}
    \item \textbf{Binary Semaphore} (\texttt{semControlDisplay}): the
          \texttt{vTaskOnOffControl} task signals this semaphore each time the relay
          state changes; \texttt{vTaskDisplay} waits for it (with a 150\,ms timeout)
          before refreshing the LCD.  This makes display updates event-driven rather
          than purely periodic, reducing unnecessary I2C traffic.
    \item \textbf{Mutex} (\texttt{lcdMutex}): protects the I2C LCD from concurrent
          write corruption.  All calls to the \texttt{LcdI2c} driver are wrapped in
          mutex take/give pairs with a 100\,ms timeout to prevent deadlocks.
\end{itemize}

% ============================================================================
% CHAPTER 3: RESOURCES AND MATERIALS
% ============================================================================
\section*{3. Resources and Materials}

\subsection*{3.1. Hardware Resources}

\begin{table}[H]
    \centering
    \caption{Hardware Components}
    \label{tab:hardware-resources}
    \begin{tabular}{|l|l|l|}
        \hline
        \textbf{Component} & \textbf{Specification} & \textbf{Quantity} \\
        \hline
        ESP32 DevKit V1 & Dual-core 240\,MHz, 520\,KB SRAM, 4\,MB Flash & 1 \\
        \hline
        DHT11 Sensor & Digital temp./humidity, single-wire, GPIO~5 & 1 \\
        \hline
        Relay Module & 5\,V relay, opto-isolated, GPIO~23 & 1 \\
        \hline
        Push Button (Joystick SW) & Active-low with pull-up, GPIO~18 & 1 \\
        \hline
        LED Indicator & 220\,$\Omega$ resistor, GPIO~2 & 1 \\
        \hline
        LCD 16$\times$2 & I2C interface, address 0x27, SDA~21 / SCL~22 & 1 \\
        \hline
        Breadboard & 400/830 tie-points & 1 \\
        \hline
        Jumper Wires & Male-to-male, male-to-female & \textasciitilde20 \\
        \hline
        USB Cable & Micro-USB for power/programming & 1 \\
        \hline
    \end{tabular}
\end{table}

\subsection*{3.2. Software Resources}

\begin{table}[H]
    \centering
    \caption{Software Tools and Frameworks}
    \label{tab:software-resources}
    \begin{tabular}{|l|l|l|}
        \hline
        \textbf{Tool / Framework} & \textbf{Version} & \textbf{Purpose} \\
        \hline
        PlatformIO & Latest & Build system, library management \\
        \hline
        FreeRTOS & 10.4.6 (bundled) & Real-time kernel \\
        \hline
        Arduino Framework ESP32 & Espressif32 6.9.0 & Hardware abstraction, GPIO \\
        \hline
        Adafruit DHT Library & 1.4.6 & DHT11 sensor driver \\
        \hline
        Adafruit Unified Sensor & 1.1.14 & Sensor abstraction base \\
        \hline
        LiquidCrystal\_I2C & 1.1.4 & 16$\times$2 I2C LCD driver \\
        \hline
        VS Code + PlatformIO ext. & Latest & IDE and debugger \\
        \hline
        Git & 2.43.0 & Version control \\
        \hline
        Arduino Serial Plotter & Built-in & Real-time data visualisation \\
        \hline
    \end{tabular}
\end{table}

% ============================================================================
% CHAPTER 4: LABORATORY TASK DESCRIPTION
% ============================================================================
\section*{4. Laboratory Task Description}

\subsection*{4.1. Problem Statement}

Design and implement a modular, FreeRTOS-based embedded application for the ESP32
that realises ON-OFF temperature control with hysteresis (Variant~A). The system
must read temperature from a DHT11 digital sensor, compare it with a configurable
Set Point, and actuate a relay accordingly. Hysteresis must be applied to prevent
frequent relay toggling near the Set Point. The Set Point must be adjustable at
runtime through both a hardware push button and serial commands. Measured
temperature, Set Point, relay state, and hysteresis value must be displayed on a
16$\times$2 I2C LCD. Real-time data must be streamed to the Arduino Serial Plotter
in the format \texttt{SetPoint:<v> Temp:<v> Output:<v>}. All functional blocks
must run as independent FreeRTOS tasks with justified recurrence periods.

\subsection*{4.2. Functional Requirements}

\begin{table}[H]
    \centering
    \caption{Functional Requirements}
    \label{tab:functional-requirements}
    \begin{tabular}{|l|l|}
        \hline
        \textbf{Requirement} & \textbf{Description} \\
        \hline
        Temperature Acquisition & Read DHT11 every 1\,s; simulate if sensor absent \\
        \hline
        ON-OFF Control with Hysteresis & Two-point relay control inside FreeRTOS task \\
        \hline
        Set Point Configuration (Button) & Short press $+1\,^{\circ}$C; long press $-1\,^{\circ}$C \\
        \hline
        Set Point Configuration (Serial) & Command \texttt{sp<val>}, range 20--60\,$^{\circ}$C \\
        \hline
        Hysteresis Configuration & Serial command \texttt{hyst<val>}, range 0.1--10\,$^{\circ}$C \\
        \hline
        Emergency Stop / Resume & Commands \texttt{stop} / \texttt{run} via serial \\
        \hline
        Status Report & Serial command \texttt{status} prints current state \\
        \hline
        LCD Display & SP, temperature, relay state, hysteresis on 16$\times$2 LCD \\
        \hline
        Serial Plotter Output & \texttt{SetPoint} / \texttt{Temp} / \texttt{Output} stream \\
        \hline
        Thread Safety & Mutex for LCD; semaphore for control$\to$display events \\
        \hline
    \end{tabular}
\end{table}

\subsection*{4.3. Non-Functional Requirements}

\begin{table}[H]
    \centering
    \caption{Non-Functional Requirements}
    \label{tab:non-functional-requirements}
    \begin{tabular}{|l|l|l|}
        \hline
        \textbf{Requirement} & \textbf{Target} & \textbf{Acceptance Criteria} \\
        \hline
        Control Period & 100\,ms & \texttt{vTaskDelayUntil} timing \\
        \hline
        Acquisition Period & 1000\,ms & DHT11 minimum sample interval \\
        \hline
        Display Period & 500\,ms & Event-driven + periodic fallback \\
        \hline
        Command Latency & $<$200\,ms & Serial input to relay response \\
        \hline
        No Relay Chattering & Zero toggles within deadband & Oscilloscope / log verification \\
        \hline
        Thread Safety & No LCD corruption & Mutex timeout protection \\
        \hline
        Modular Design & Clear HAL/App separation & Code review \\
        \hline
        Code Portability & Reusable drivers & Build with bare-metal HAL \\
        \hline
    \end{tabular}
\end{table}

\subsection*{4.4. Two-Stage Methodology}

\textbf{Stage 1 – Basic ON-OFF Control}

Initialise FreeRTOS, create the three tasks, and verify basic relay toggling in
response to a temperature threshold. Implement the DHT11 driver and confirm that raw
temperature readings are printed to the serial monitor. Validate that the relay
GPIO drives the LED in synchrony with the control output.

\textbf{Stage 2 – Full System with Hysteresis and Interactive Configuration}

Add the two-point hysteresis logic to eliminate chattering. Integrate the push
button with short/long press detection for Set Point adjustment. Implement the
serial command parser (\texttt{sp}, \texttt{hyst}, \texttt{status}, \texttt{stop},
\texttt{run}). Add LCD display via the I2C driver and protect access with a mutex.
Connect the binary semaphore for event-driven display refresh. Validate all
requirements and test edge cases (temperature at boundary, emergency stop, sensor
disconnected).

% ============================================================================
% CHAPTER 5: TEST SCENARIOS AND VALIDATION CRITERIA
% ============================================================================
\section*{5. Test Scenarios and Validation Criteria}

\subsection*{5.1. Test Scenarios}

\textbf{Test Scenario 1: System Initialisation} -- verifies that all FreeRTOS
tasks start successfully, hardware components initialise without error, and the LCD
shows \texttt{"Lab 5.1 Ready"} followed by \texttt{"Temp ON-OFF Ctrl"} on power-up.

\textbf{Test Scenario 2: Hysteresis Deadband Verification} -- sets a known Set
Point (e.g.\ 30\,$^{\circ}$C) and hysteresis (2\,$^{\circ}$C) via serial commands.
Verifies that no relay toggling occurs while the temperature remains inside
$[28,32]\,^{\circ}$C, and that exactly one toggle occurs when each boundary is
crossed.

\textbf{Test Scenario 3: Button Set Point Adjustment} -- performs short and long
button presses and verifies that the Set Point increments by $1\,^{\circ}$C on a
short press and decrements by $1\,^{\circ}$C on a long press. Verifies that the
Set Point is clamped to $[20, 60]\,^{\circ}$C.

\textbf{Test Scenario 4: Serial Command Interface} -- sends commands
\texttt{sp35}, \texttt{hyst2}, \texttt{status}, \texttt{stop}, and \texttt{run}
via the serial monitor. Verifies correct parsing, acknowledgement messages, and
relay behaviour for each command.

\textbf{Test Scenario 5: Simulation Mode (No Sensor)} -- disconnects the DHT11
and verifies that the system falls back to simulation. Confirms that temperature
rises when the relay is ON (+0.2\,$^{\circ}$C/cycle) and falls when OFF
(-0.1\,$^{\circ}$C/cycle), and that the system oscillates stably around the Set
Point.

\textbf{Test Scenario 6: LCD and Serial Plotter Output} -- confirms that the LCD
line~1 shows the format \texttt{"SP:xx.x T:xx.xC"} and line~2 shows
\texttt{"Rel:ON/OFF H:x.xC"}, and that the serial output follows the Arduino
Serial Plotter format with labelled fields.

\subsection*{5.2. Validation Criteria}

\begin{table}[H]
    \centering
    \caption{Validation Criteria}
    \label{tab:validation-criteria}
    \begin{tabular}{|l|l|l|}
        \hline
        \textbf{Criterion} & \textbf{Requirement} & \textbf{Measurement Method} \\
        \hline
        No chattering in deadband & 0 toggles while $|T - SP| < H$ & Serial log inspection \\
        \hline
        Control period & 100\,ms $\pm$5\,ms & FreeRTOS tick counter \\
        \hline
        Display period & 500\,ms $\pm$50\,ms & Stopwatch / serial timestamps \\
        \hline
        Command latency & $<$200\,ms & Serial command to relay GPIO \\
        \hline
        Button responsiveness & $<$150\,ms from release & Manual test \\
        \hline
        Set Point clamping & Stays in [20, 60]\,$^{\circ}$C & Boundary value test \\
        \hline
        CPU utilisation & $<$50\% & FreeRTOS runtime stats \\
        \hline
        Memory usage & $<$50\,KB & Heap analysis \\
        \hline
    \end{tabular}
\end{table}

% ============================================================================
% CHAPTER 6: ARCHITECTURAL DESIGN AND BEHAVIORAL MODELING
% ============================================================================
\section*{6. Architectural Design and Behavioral Modeling}

\subsection*{6.1. High-Level Architecture}

The system is organised into four horizontal layers as shown in
Figure~\ref{fig:architecture}.

\begin{figure}[H]
\centering
\begin{tikzpicture}[
    layer/.style={draw, rounded corners=4pt, minimum width=13cm,
                  minimum height=1.2cm, align=center, font=\small},
    box/.style={draw, rounded corners=3pt, fill=blue!15,
                minimum width=2.8cm, minimum height=0.9cm, align=center,
                font=\footnotesize},
    greybox/.style={draw, rounded corners=3pt, fill=gray!15,
                    minimum width=2.8cm, minimum height=0.9cm, align=center,
                    font=\footnotesize},
]
% Layer 4 – Hardware
\node[layer, fill=red!10]   (hw)   at (0,0)   {};
\node[font=\bfseries\footnotesize, anchor=west] at (-6.3, 0) {Hardware Layer};
\node[greybox] at (-3.8, 0) {DHT11\\Sensor};
\node[greybox] at (-0.9, 0) {Relay\\Module};
\node[greybox] at (1.9, 0)  {LCD 16x2\\I2C};
\node[greybox] at (4.7, 0)  {Button\\+ LED};

% Layer 3 – HAL
\node[layer, fill=orange!10] (hal) at (0,1.6) {};
\node[font=\bfseries\footnotesize, anchor=west] at (-6.3, 1.6) {HAL (Drivers)};
\node[box] at (-3.8, 1.6) {Dht11Sensor};
\node[box] at (-0.9, 1.6) {Actuator};
\node[box] at (1.9, 1.6)  {LcdI2c};
\node[box] at (4.7, 1.6)  {Joystick\\Led};

% Layer 2 – Kernel Primitives
\node[layer, fill=yellow!10] (kp) at (0,3.2) {};
\node[font=\bfseries\footnotesize, anchor=west] at (-6.3, 3.2) {Kernel Primitives};
\node[box] at (-3.0, 3.2) {Mutex\\(lcdMutex)};
\node[box] at ( 0.0, 3.2) {BinarySemaphore\\(semControlDisplay)};
\node[box] at ( 3.2, 3.2) {Task\\(createTask)};

% Layer 1 – Application
\node[layer, fill=green!10] (app) at (0,4.8) {};
\node[font=\bfseries\footnotesize, anchor=west] at (-6.3, 4.8) {Application Layer};
\node[box, fill=green!25] at (-3.8, 4.8) {vTaskAcquisition\\(1000 ms, P3)};
\node[box, fill=green!25] at ( 0.0, 4.8) {vTaskOnOffControl\\(100 ms, P3)};
\node[box, fill=green!25] at ( 3.8, 4.8) {vTaskDisplay\\(500 ms, P2)};

% User input label
\node[font=\bfseries\footnotesize] at (0, 6.2) {User Input: Button / Serial Commands};
\draw[->, thick] (0,6.0) -- (0,5.35);
\end{tikzpicture}
\caption{High-Level Architecture -- ON-OFF Temperature Control System}
\label{fig:architecture}
\end{figure}

\subsubsection*{Hardware Layer}
Contains the physical components: DHT11 sensor, relay module, 16$\times$2 I2C LCD,
push button, and status LED.

\subsubsection*{Hardware Abstraction Layer (HAL)}
Device drivers wrap GPIO, I2C, and sensor protocol specifics, exposing clean C++
class interfaces to upper layers.

\subsubsection*{Kernel Primitives Layer}
FreeRTOS wrapper classes simplify RTOS API usage: \texttt{Mutex},
\texttt{BinarySemaphore}, and \texttt{createTask} hide raw FreeRTOS handles and
provide RAII-style resource management.

\subsubsection*{Application Layer}
Three FreeRTOS tasks implement the system behaviour: data acquisition, ON-OFF
control with hysteresis, and display updates.

\subsection*{6.2. Component Diagram}

% Layout: left-to-right pipeline, data flows west→east.
% acq and ctrl feed SharedData from opposite diagonals (symmetric,
% no crossing).  The semaphore arc stays above y = −1.5, below the
% horizontal chain, so it never intersects any other arrow.

\begin{figure}[H]
\centering
\resizebox{\linewidth}{!}{%
\begin{tikzpicture}[
    comp/.style={draw, rounded corners=2pt, fill=blue!15,
                 minimum width=2.4cm, minimum height=0.8cm,
                 align=center, font=\scriptsize},
    drv/.style={draw, rounded corners=2pt, fill=gray!20,
                minimum width=1.8cm, minimum height=0.7cm,
                align=center, font=\scriptsize},
    shared/.style={draw, rounded corners=2pt, fill=yellow!30,
                   minimum width=2.0cm, minimum height=1.1cm,
                   align=center, font=\scriptsize},
    arr/.style={->, semithick},
    sem/.style={->, semithick, dashed, blue!70!black}
]

%% ── Input drivers (far left) ──────────────────────────────
\node[drv]  (dht) at (-5.5,  1.5) {Dht11Sensor};
\node[drv]  (btn) at (-5.5, -1.5) {Button};

%% ── FreeRTOS Tasks (centre-left) ─────────────────────────
\node[comp] (acq)  at (-2.5,  1.5) {vTaskAcquisition\\1000\,ms · P3};
\node[comp] (ctrl) at (-2.5, -1.5) {vTaskOnOffControl\\100\,ms · P3};

%% ── SharedData hub (centre) ──────────────────────────────
\node[shared] (sd) at (0.5, 0.0) {SharedData\\setpoint / temp\\relay / hyst};

%% ── Display chain (right) ────────────────────────────────
\node[comp] (disp) at (3.5, 0.0) {vTaskDisplay\\500\,ms · P2};
\node[drv]  (lcd)  at (6.2, 0.0) {LcdI2c\\mutex};

%% ── Output drivers (bottom) ──────────────────────────────
\node[drv] (act) at (-1.0, -3.0) {Actuator};
\node[drv] (led) at ( 1.0, -3.0) {Led};

%% ── Arrows ───────────────────────────────────────────────

% Horizontal: drivers → tasks (same y, no crossing possible)
\draw[arr] (dht) -- node[above,font=\tiny]{reads} (acq);
\draw[arr] (btn) -- node[above,font=\tiny]{press} (ctrl);

% Symmetric diagonals into SharedData.
% acq→sd goes ↘ and ctrl→sd goes ↗; they share the same
% x-parametric so they converge to sd without ever crossing.
\draw[arr] (acq)  -- node[right,font=\tiny]{temp} (sd);
\draw[arr] (ctrl) -- node[right,font=\tiny]{SP/relay} (sd);

% Horizontal chain: SharedData → Display → LCD
\draw[arr] (sd)   -- node[above,font=\tiny]{reads} (disp);
\draw[arr] (disp) -- node[above,font=\tiny]{I2C}   (lcd);

% Output GPIO controls (go downward from ctrl, open space)
\draw[arr] (ctrl) -- node[left, font=\tiny]{GPIO} (act);
\draw[arr] (ctrl) -- node[right,font=\tiny]{GPIO} (led);

% Semaphore ctrl → disp.
% bend right=12 bows the arc gently south; the apex passes
% through ≈(0.7, −0.75), well below sd at (0.5, 0) and above
% the actuator drivers at y = −3.
\draw[sem] (ctrl) to[bend right=12]
    node[below,font=\tiny]{semaphore} (disp);

\end{tikzpicture}%
}
\caption{Component Diagram -- ON-OFF Temperature Control System}
\label{fig:component}
\end{figure}

\subsection*{6.3. Class Diagram}

Key classes are shown below:
\begin{itemize}
    \item \textbf{Dht11Sensor}: wraps Adafruit DHT; exposes
          \texttt{readTemperature()}, \texttt{readHumidity()},
          \texttt{isPresent()}.
    \item \textbf{Actuator}: binary relay control; exposes \texttt{turnOn()},
          \texttt{turnOff()}, \texttt{isActive()}, \texttt{getLastToggleTime()}.
    \item \textbf{LcdI2c}: I2C LCD driver; exposes \texttt{print()},
          \texttt{setCursor()}, \texttt{clear()}, \texttt{printf()}.
    \item \textbf{Joystick}: button with press-duration detection;
          \texttt{scan()}, \texttt{getPressDuration()}.
    \item \textbf{Led}: simple GPIO toggle wrapper.
    \item \textbf{Mutex} / \textbf{BinarySemaphore}: FreeRTOS API wrappers with
          RAII-style \texttt{init()}, \texttt{take()}, \texttt{give()}.
\end{itemize}

\subsection*{6.4. Task Activity Diagrams}

\subsubsection*{vTaskAcquisition Activity (1000\,ms, Priority~3)}

\begin{enumerate}
    \item Initialise \texttt{xLastWakeTime} from current tick.
    \item Enter infinite loop.
    \item Block until 1000\,ms elapsed via \texttt{vTaskDelayUntil()}.
    \item Call \texttt{dht11Sensor->readTemperature()}.
    \item If reading valid ($>-100\,^{\circ}$C): update \texttt{sharedData.temperature}.
    \item Else (sensor absent): simulate heating ($+0.2\,^{\circ}$C when relay ON)
          or cooling ($-0.1\,^{\circ}$C when relay OFF), clamped to
          $[15, 70]\,^{\circ}$C.
    \item Repeat.
\end{enumerate}

\subsubsection*{vTaskOnOffControl Activity (100\,ms, Priority~3)}

\begin{enumerate}
    \item Initialise \texttt{xLastWakeTime}; set \texttt{emergencyStop = false}.
    \item Enter infinite loop.
    \item Block until 100\,ms elapsed.
    \item Process serial command if \texttt{serial\_command\_received}: handle
          \texttt{stop}, \texttt{run}, \texttt{status}, \texttt{hyst<v>},
          \texttt{sp<v>}.
    \item Scan joystick button via \texttt{joystick->scan()}; read press duration.
    \item If short press ($<$500\,ms): \texttt{setpoint += 1.0°C} (clamped).
    \item If long press ($\ge$500\,ms): \texttt{setpoint -= 1.0°C} (clamped).
    \item If \texttt{emergencyStop}: skip control, continue.
    \item Apply two-point hysteresis: compute \texttt{newOn} from
          \texttt{temp}, \texttt{setpoint}, \texttt{hysteresis}.
    \item If state changed: update relay and LED; print log.
    \item Signal \texttt{semControlDisplay.give()}.
    \item Repeat.
\end{enumerate}

\subsubsection*{vTaskDisplay Activity (500\,ms, Priority~2)}

\begin{enumerate}
    \item Initialise \texttt{xLastWakeTime}.
    \item Enter infinite loop.
    \item Block until 500\,ms elapsed.
    \item Wait for \texttt{semControlDisplay.take(150\,ms)} -- event or timeout.
    \item Format LCD line 1: \texttt{"SP:xx.x T:xx.xC"}.
    \item Format LCD line 2: \texttt{"Rel:ON/OFF H:x.xC"}.
    \item Call \texttt{updateLCD(line1, line2)} (mutex-protected).
    \item Print Serial Plotter line:
          \texttt{"SetPoint:xx.x Temp:xx.x Output:0/1"}.
    \item Repeat.
\end{enumerate}

\subsection*{6.5. Data Flow}

\begin{enumerate}
    \item DHT11 sensor $\to$ \texttt{vTaskAcquisition} $\to$
          \texttt{sharedData.temperature}.
    \item Button / Serial $\to$ \texttt{vTaskOnOffControl} $\to$
          \texttt{sharedData.setpoint}, \texttt{sharedData.hysteresis}.
    \item Hysteresis logic $\to$ \texttt{sharedData.relayOn} $\to$ Relay GPIO.
    \item \texttt{vTaskOnOffControl} $\xrightarrow{\text{semaphore}}$
          \texttt{vTaskDisplay} $\to$ LCD / Serial Plotter.
\end{enumerate}

\subsection*{6.6. State Machine}

\begin{figure}[H]
\centering
\begin{tikzpicture}[
    >=stealth,
    state/.style={draw, circle, minimum size=2.8cm, align=center, font=\small},
    initial/.style={state, fill=green!20},
    stop/.style={state, fill=red!20},
]
\node[initial] (heat) at (-3.5, 0) {HEATING\\(Relay ON)};
\node[state, fill=blue!10] (cool) at ( 3.5, 0) {COOLING\\(Relay OFF)};
\node[stop]  (estop) at ( 0, -3.0) {EMERGENCY\\STOP};

% Normal transitions
\draw[->, thick, bend left=30]
    (heat) to node[above, font=\footnotesize]
    {$T > SP + H$} (cool);

\draw[->, thick, bend left=30]
    (cool) to node[below, font=\footnotesize]
    {$T < SP - H$} (heat);

% Self-loops (deadband)
\draw[->, thick, loop above] (heat)
    to node[above, font=\footnotesize]{deadband: hold} (heat);
\draw[->, thick, loop above] (cool)
    to node[above, font=\footnotesize]{deadband: hold} (cool);

% Emergency stop
\draw[->, thick, dashed, bend right=20]
    (heat) to node[left, font=\footnotesize]{\texttt{stop}} (estop);
\draw[->, thick, dashed, bend left=20]
    (cool) to node[right, font=\footnotesize]{\texttt{stop}} (estop);
\draw[->, thick, dashed]
    (estop) to node[right, font=\footnotesize]{\texttt{run}} (cool);
\end{tikzpicture}
\caption{State Machine -- ON-OFF Control with Hysteresis}
\label{fig:state-machine}
\end{figure}

The system has two operational states plus an emergency stop:
\begin{itemize}
    \item \textbf{HEATING} (Relay ON): entered when $T < SP - H$.
          Maintained while $T \le SP + H$ (deadband).
    \item \textbf{COOLING} (Relay OFF): entered when $T > SP + H$.
          Maintained while $T \ge SP - H$ (deadband).
    \item \textbf{EMERGENCY STOP}: relay forced OFF; control loop suspended.
          Resumed on \texttt{run} command.
\end{itemize}

\subsection*{6.7. Task Priorities}

\begin{table}[H]
    \centering
    \caption{FreeRTOS Task Priority Assignment}
    \label{tab:task-priorities}
    \begin{tabular}{|l|c|c|l|}
        \hline
        \textbf{Task} & \textbf{Priority} & \textbf{Period} & \textbf{Rationale} \\
        \hline
        \texttt{vTaskAcquisition} & P3 & 1000\,ms & DHT11 conversion $\ge$1\,s \\
        \hline
        \texttt{vTaskOnOffControl} & P3 & 100\,ms & Responsive relay switching \\
        \hline
        \texttt{vTaskDisplay} & P2 & 500\,ms & Lower priority; can be delayed \\
        \hline
    \end{tabular}
\end{table}

\subsection*{6.8. Timing Diagram}

\begin{figure}[H]
\centering
\begin{tikzpicture}[x=0.4cm, y=0.6cm]
    % Axis
    \draw[->] (0,0) -- (26,0) node[right]{\footnotesize time (ms)};
    \foreach \x/\label in {0/0, 5/500, 10/1000, 15/1500, 20/2000, 25/2500}
        \draw (\x,0.1) -- (\x,-0.1) node[below]{\footnotesize\label};

    % vTaskAcquisition (1000ms) -- row y=4
    \node[anchor=east, font=\footnotesize] at (0,4.5) {Acquisition};
    \foreach \s in {0, 10, 20}
        \filldraw[fill=green!40, draw=green!60!black]
            (\s,4) rectangle (\s+0.4, 5);

    % vTaskOnOffControl (100ms) -- row y=2.5
    \node[anchor=east, font=\footnotesize] at (0,3.0) {OnOffCtrl};
    \foreach \s in {0,1,2,3,4,5,6,7,8,9,10,11,12,13,14,15,16,17,18,19,20,21,22,23,24}
        \filldraw[fill=blue!40, draw=blue!60!black]
            (\s,2.5) rectangle (\s+0.3, 3.5);

    % vTaskDisplay (500ms) -- row y=1
    \node[anchor=east, font=\footnotesize] at (0,1.5) {Display};
    \foreach \s in {0, 5, 10, 15, 20, 25}
        \filldraw[fill=orange!40, draw=orange!60!black]
            (\s,1) rectangle (\s+0.6, 2);

    % Semaphore signal markers
    \foreach \x in {1,3,7,11,13,19}
        \draw[->, dashed, thick, blue] (\x+0.15,2.5) -- (\x+0.15, 2.1);
\end{tikzpicture}
\caption{Timing Diagram -- Task Execution and Semaphore Signals}
\label{fig:timing}
\end{figure}

Timing characteristics:
\begin{itemize}
    \item \texttt{vTaskAcquisition}: 1000\,ms period, $\approx$2\,ms execution.
    \item \texttt{vTaskOnOffControl}: 100\,ms period, $\approx$1\,ms execution.
    \item \texttt{vTaskDisplay}: 500\,ms period, $\approx$15\,ms execution (I2C).
    \item Relay switch latency: $<$1\,ms from threshold crossing to GPIO change.
    \item Dashed arrows: semaphore signals from control to display task.
\end{itemize}

% ============================================================================
% CHAPTER 7: ELECTRICAL DESIGN AND SIMULATION
% ============================================================================
\section*{7. Electrical Design and Simulation}

\subsection*{7.0. Circuit Wiring Sketch}

Figure~\ref{fig:electrical-sketch} shows the complete wiring between the ESP32 and
all peripheral modules.  Blue lines are signal wires, red/orange lines are power
supply connections, and black lines are ground connections.

% ── Electrical wiring sketch ────────────────────────────────────────────────
% Design rule: every signal wire is purely HORIZONTAL (one unique y per GPIO).
% Power / GND are shown as floating flag labels attached directly to each
% component – NO long vertical wires, so there are ZERO wire crossings.
% ─────────────────────────────────────────────────────────────────────────────

\begin{figure}[H]
\centering
\resizebox{\linewidth}{!}{%
\begin{tikzpicture}[
  sig/.style={draw, line width=1pt, blue!65!black},
  wire/.style={draw, line width=0.9pt},
  % Floating power-flag label styles
  f3v/.style={draw=red!70!black, fill=red!10, rounded corners=1.5pt,
              font=\tiny\bfseries, inner sep=2pt, text=red!80!black},
  f5v/.style={draw=orange!65!black, fill=orange!10, rounded corners=1.5pt,
              font=\tiny\bfseries, inner sep=2pt, text=orange!80!black},
  fgnd/.style={draw=black!55, fill=gray!10, rounded corners=1.5pt,
               font=\tiny\bfseries, inner sep=2pt},
]

%% ── ESP32 board ──────────────────────────────────────────────────────────
\draw[fill=green!12, line width=1pt, rounded corners=5pt]
    (0.2, 0.3) rectangle (3.2, 9.1);
\node[font=\small\bfseries, rotate=90] at (1.7, 4.7) {ESP32 DevKit V1};

% Six GPIO stubs on the right edge (each at its own y – never cross)
\foreach \y/\lbl in {
    8.0/{GPIO 5},
    6.5/{GPIO 23},
    5.0/{GPIO 18},
    3.5/{GPIO 2},
    2.3/{SDA 21},
    1.0/{SCL 22}}{
  \draw[sig] (3.2,\y) -- (3.6,\y);
  \node[anchor=east, font=\tiny] at (3.2,\y) {\lbl};
}

% Power-flag annotations on the left edge of the ESP32
\node[f3v,  anchor=east] at (0.2, 8.6) {+3.3V};
\node[f5v,  anchor=east] at (0.2, 7.8) {+5V};
\node[fgnd, anchor=east] at (0.2, 1.2) {GND};

%% ── DHT11 Sensor ─────────────────────────────────────────────────────────
% Box (center y = 8.0, same as GPIO 5 stub)
\draw[fill=blue!12, line width=0.9pt, rounded corners=3pt]
    (5.2, 7.55) rectangle (8.0, 8.45);
\node[font=\scriptsize\bfseries] at (6.6, 8.12) {DHT11};
\node[font=\tiny]                at (6.6, 7.78) {Temp/Humidity};

% Signal wire: GPIO5 → DHT11 (horizontal, y=8.0)
\draw[sig] (3.6,8.0) -- (5.2,8.0);
\node[above, font=\tiny, blue!65!black] at (4.4,8.0) {DATA};

% Pull-up resistor: short vertical branch upward from junction at x=4.3
% (stays above y=8.0; nothing else exists above this level)
\filldraw[black] (4.3,8.0) circle (0.06);          % junction dot
\draw[wire] (4.3,8.0) -- (4.3,8.55);              % wire up to resistor
\draw[wire] (4.08,8.55) rectangle (4.52,8.95);     % resistor body
\node[font=\tiny] at (4.3,8.75) {10k$\Omega$};
\draw[wire] (4.3,8.95) -- (4.3,9.2);              % wire to flag
\node[f3v] at (4.3,9.4) {+3.3V};                  % power flag (top)

% Power flags attached to right edge of DHT11 box
\draw[wire, thin] (8.0,8.25) -- (8.2,8.25);
\node[f3v,  anchor=west] at (8.2,8.25) {+3.3V};
\draw[wire, thin] (8.0,7.75) -- (8.2,7.75);
\node[fgnd, anchor=west] at (8.2,7.75) {GND};

%% ── Relay Module ─────────────────────────────────────────────────────────
% Box (center y = 6.5, same as GPIO 23 stub)
\draw[fill=orange!15, line width=0.9pt, rounded corners=3pt]
    (5.2, 6.05) rectangle (8.0, 6.95);
\node[font=\scriptsize\bfseries] at (6.3, 6.60) {Relay Module};
\node[font=\tiny]                at (6.3, 6.22) {5V opto-isolated};

% Signal wire: GPIO23 → Relay (horizontal, y=6.5)
\draw[sig] (3.6,6.5) -- (5.2,6.5);
\node[above, font=\tiny, blue!65!black] at (4.4,6.5) {IN};

% Coil symbol inside box (right half)
\draw[wire, rounded corners=0.5pt]
    (6.9,6.25)--(7.0,6.50)--(7.1,6.32)--(7.2,6.58)
    --(7.3,6.38)--(7.4,6.62)--(7.5,6.48)--(7.6,6.25);
\node[font=\tiny] at (7.25,6.80) {coil};

% Load arrow on right
\draw[->, line width=1.1pt, red!60!black] (8.0,6.5) -- (8.7,6.5);
\node[font=\tiny, anchor=west, red!60!black] at (8.7,6.5) {Load};

% Power flags
\draw[wire, thin] (8.0,6.78) -- (8.2,6.78);
\node[f5v,  anchor=west] at (8.2,6.78) {+5V};
\draw[wire, thin] (8.0,6.22) -- (8.2,6.22);
\node[fgnd, anchor=west] at (8.2,6.22) {GND};

%% ── Push Button (normally-open, active-low) ──────────────────────────────
% Signal wire: GPIO18 → left terminal (horizontal, y=5.0)
\draw[sig] (3.6,5.0) -- (5.5,5.0);
\node[above, font=\tiny, blue!65!black] at (4.4,5.0) {SW};

% Switch symbol at x≈5.5–6.2
\filldraw[draw=blue!65!black, fill=white, line width=0.8pt] (5.5,5.0)  circle (0.09);
\draw[sig, line width=1pt] (5.59,5.08) -- (6.11,5.28);   % open arm
\filldraw[draw=blue!65!black, fill=white, line width=0.8pt] (6.21,5.0) circle (0.09);
\node[font=\scriptsize\bfseries] at (5.85,5.52) {Button};

% GND flag directly on right terminal (no wire going down)
\draw[wire, thin] (6.30,5.0) -- (6.50,5.0);
\node[fgnd, anchor=west] at (6.5,5.0) {GND};

%% ── LED + 220 Ω resistor ─────────────────────────────────────────────────
% Signal wire: GPIO2 → resistor (horizontal, y=3.5)
\draw[sig] (3.6,3.5) -- (4.5,3.5);

% Resistor body
\draw[wire] (4.5,3.30) rectangle (5.1,3.70);
\node[font=\tiny] at (4.8,3.50) {220$\Omega$};

% Wire resistor → LED anode
\draw[sig] (5.1,3.5) -- (5.55,3.5);

% LED symbol: filled triangle + cathode bar
\fill[yellow!55!orange] (5.55,3.22) -- (5.55,3.78) -- (6.05,3.50) -- cycle;
\draw[wire]             (5.55,3.22) -- (5.55,3.78) -- (6.05,3.50) -- cycle;
\draw[wire, line width=1.3pt] (6.05,3.22) -- (6.05,3.78);

% Light rays
\draw[yellow!40!orange, ->, line width=0.7pt] (5.9,3.85) -- (6.15,4.15);
\draw[yellow!40!orange, ->, line width=0.7pt] (6.05,3.80) -- (6.35,4.08);
\node[font=\tiny] at (5.8,3.10) {LED};

% GND flag on cathode (no long wire down)
\draw[wire, thin] (6.05,3.5) -- (6.30,3.5);
\node[fgnd, anchor=west] at (6.30,3.5) {GND};

%% ── LCD 16×2 I2C ─────────────────────────────────────────────────────────
% Box spans y=0.6–2.8; SDA enters at y=2.3, SCL at y=1.0
\draw[fill=cyan!13, line width=0.9pt, rounded corners=3pt]
    (5.2, 0.6) rectangle (8.0, 2.8);
\node[font=\scriptsize\bfseries] at (6.6, 2.52) {LCD 16$\times$2 I2C};
\node[font=\tiny]                at (6.6, 2.20) {I2C Addr: 0x27};

% Simulated character-cell grid (two rows inside the lower half of box)
\foreach \xi in {0,...,5}{
  \draw[gray!45, thin]
      ({5.35+\xi*0.44},0.72) rectangle ({5.35+\xi*0.44+0.36},1.05);
  \draw[gray!45, thin]
      ({5.35+\xi*0.44},1.13) rectangle ({5.35+\xi*0.44+0.36},1.46);
}

% SDA wire: GPIO21 → LCD left edge (horizontal, y=2.3)
\draw[sig] (3.6,2.3) -- (5.2,2.3);
\node[above, font=\tiny, blue!65!black] at (4.4,2.3) {SDA};

% SCL wire: GPIO22 → LCD left edge (horizontal, y=1.0)
\draw[sig] (3.6,1.0) -- (5.2,1.0);
\node[above, font=\tiny, blue!65!black] at (4.4,1.0) {SCL};

% Power flags on right edge of LCD
\draw[wire, thin] (8.0,2.60) -- (8.2,2.60);
\node[f5v,  anchor=west] at (8.2,2.60) {+5V};
\draw[wire, thin] (8.0,0.80) -- (8.2,0.80);
\node[fgnd, anchor=west] at (8.2,0.80) {GND};

\end{tikzpicture}%
}
\caption{Electrical Circuit Wiring Sketch -- ON-OFF Temperature Control System}
\label{fig:electrical-sketch}
\end{figure}

\subsection*{7.1. Hardware Connections}

\begin{table}[H]
    \centering
    \caption{GPIO Pin Assignments}
    \label{tab:pin-assignments}
    \begin{tabular}{|l|l|l|}
        \hline
        \textbf{GPIO Pin} & \textbf{Component} & \textbf{Function} \\
        \hline
        GPIO~5 & DHT11 DATA & 1-wire digital temperature/humidity \\
        \hline
        GPIO~23 & Relay IN & Binary actuator control (HIGH = ON) \\
        \hline
        GPIO~18 & Button (Joystick SW) & Digital input, INPUT\_PULLUP \\
        \hline
        GPIO~2 & LED & Status indicator (mirrors relay state) \\
        \hline
        GPIO~21 (SDA) & LCD I2C & I2C data line \\
        \hline
        GPIO~22 (SCL) & LCD I2C & I2C clock line \\
        \hline
        3.3V & DHT11 VCC & Sensor power supply \\
        \hline
        5V & Relay VCC, LCD VCC & Module power supply \\
        \hline
        GND & All modules & Common ground \\
        \hline
    \end{tabular}
\end{table}

\subsection*{7.2. DHT11 Sensor Circuit}

The DHT11 DATA pin is connected to GPIO~5 through a $10\,\text{k}\Omega$ pull-up
resistor to 3.3\,V.  The sensor is powered from the ESP32 3.3\,V rail.  The single
data wire carries both clock and data during the proprietary single-wire protocol
handled by the Adafruit DHT library.

\subsection*{7.3. Relay Circuit}

The relay module integrates an opto-coupler and transistor driver.  GPIO~23 is
connected to the relay \texttt{IN} pin.  Driving the pin HIGH energises the relay
coil and closes the NO contact; driving it LOW de-energises the coil.  A flyback
diode on the module suppresses inductive voltage spikes.  The relay coil is powered
from the 5\,V rail to avoid loading the ESP32 3.3\,V LDO.

\subsection*{7.4. LCD I2C Circuit}

The 16$\times$2 LCD uses an I2C backpack at address 0x27.  The SDA and SCL lines
are connected to GPIO~21 and GPIO~22 respectively.  Internal 4.7\,k$\Omega$
pull-up resistors on the backpack hold the bus HIGH when idle.  The LCD is powered
from 5\,V; the I2C signal levels are compatible with the ESP32 3.3\,V I/O due to
the backpack level shifter.

\subsection*{7.5. Power Supply Considerations}

The USB 5\,V rail powers the relay module and LCD.  The ESP32 onboard LDO provides
3.3\,V for the microcontroller and DHT11.  Estimated current budget:
\begin{itemize}
    \item ESP32 active: 80--150\,mA
    \item DHT11: $<$2\,mA
    \item Relay coil: 50--70\,mA (when ON)
    \item LCD + backlight: 20--30\,mA
    \item LED: 6\,mA
    \item \textbf{Total}: $\approx$160--260\,mA (USB-powered, within 500\,mA limit)
\end{itemize}

% ============================================================================
% CHAPTER 8: SOFTWARE IMPLEMENTATION
% ============================================================================
\section*{8. Software Implementation}

\subsection*{8.1. Project Structure}

\begin{verbatim}
ES/
|-- src/
|   |-- main.cpp
|   `-- modules/
|       |-- freertos_app/
|       |   |-- freertos_app.h/cpp      # Public setup/loop interface
|       |   |-- init/init.h/cpp         # Hardware and task initialisation
|       |   |-- state/state.h/cpp       # Shared data, constants, HW objects
|       |   |-- sync/sync.h/cpp         # Sync primitives & updateLCD()
|       |   `-- tasks/tasks.h/cpp       # Three FreeRTOS tasks
|       |-- kernel_primitives/
|       |   |-- task/task.h/cpp         # createTask / delayMs wrappers
|       |   |-- mutex/mutex.h/cpp
|       |   `-- semaphore/binary_semaphore.h/cpp
|       |-- dht11/dht11.h/cpp           # Adafruit DHT wrapper
|       |-- actuator/actuator.h/cpp     # Relay GPIO driver
|       |-- lcd/lcd.h/cpp               # LiquidCrystal_I2C wrapper
|       |-- joystick/joystick.h/cpp     # Button press-duration detector
|       |-- led/led.h/cpp               # LED GPIO wrapper
|       `-- serial_stdio/serial_stdio.h/cpp  # Non-blocking serial parser
|-- platformio.ini
`-- lib/lab5.1/
\end{verbatim}

\subsection*{8.2. Shared Data Structure}

\begin{lstlisting}[language=C++, caption=Shared Data Structure (state.h)]
struct SharedData {
    // Serial command interface
    bool    serial_command_received;
    char    serial_command_buffer[16];
    uint8_t serial_command_index;

    // Lab 5.1 ON-OFF Temperature Control with Hysteresis
    float setpoint;     // target temperature [degC]
    float temperature;  // measured temperature [degC]
    bool  relayOn;      // current relay state
    float hysteresis;   // deadband [degC]
};
\end{lstlisting}

\subsection*{8.3. Configuration Constants}

\begin{lstlisting}[language=C++, caption=Configuration Constants (state.cpp)]
// GPIO pins
const uint8_t DHT11_PIN      = 5;
const uint8_t RELAY_PIN      = 23;
const uint8_t JOYSTICK_SW_PIN= 18;
const uint8_t LED_PIN        = 2;
const uint8_t LCD_SDA_PIN    = 21;
const uint8_t LCD_SCL_PIN    = 22;

// Task configuration
const uint32_t    TASK_STACK_SIZE           = 4096;
const UBaseType_t TASK_PRIORITY_DISPLAY     = tskIDLE_PRIORITY + 2;
const UBaseType_t TASK_PRIORITY_ACQUISITION = tskIDLE_PRIORITY + 3;
const UBaseType_t TASK_PRIORITY_CONTROL     = tskIDLE_PRIORITY + 3;

// Task periods
const uint32_t ACQUISITION_PERIOD_MS = 1000; // DHT11 min interval
const uint32_t CONTROL_PERIOD_MS     = 100;
const uint32_t DISPLAY_PERIOD_MS     = 500;

// Control constants
const float DEFAULT_SETPOINT   = 23.0f;  // degC
const float DEFAULT_HYSTERESIS =  1.0f;  // degC (total deadband = 2 degC)
const float SETPOINT_MIN       = 20.0f;
const float SETPOINT_MAX       = 60.0f;
const float SETPOINT_STEP      =  1.0f;
\end{lstlisting}

\subsection*{8.4. Acquisition Task}

\begin{lstlisting}[language=C++, caption=vTaskAcquisition (tasks.cpp)]
void vTaskAcquisition(void *pvParameters) {
    (void)pvParameters;
    TickType_t xLastWakeTime = xTaskGetTickCount();
    const TickType_t xFrequency = pdMS_TO_TICKS(ACQUISITION_PERIOD_MS);

    for (;;) {
        vTaskDelayUntil(&xLastWakeTime, xFrequency);

        if (dht11Sensor != nullptr) {
            float t = dht11Sensor->readTemperature();
            if (t > -100.0f) {
                sharedData.temperature = t;
            } else {
                // Simulate heating/cooling when sensor is absent
                float &temp = sharedData.temperature;
                if (sharedData.relayOn) temp += 0.2f;
                else                   temp -= 0.1f;
                if (temp < 15.0f) temp = 15.0f;
                if (temp > 70.0f) temp = 70.0f;
            }
        }
    }
}
\end{lstlisting}

\subsection*{8.5. ON-OFF Control Task with Hysteresis}

\begin{lstlisting}[language=C++, caption=vTaskOnOffControl -- hysteresis core (tasks.cpp)]
void vTaskOnOffControl(void *pvParameters) {
    (void)pvParameters;
    TickType_t xLastWakeTime = xTaskGetTickCount();
    const TickType_t xFrequency = pdMS_TO_TICKS(CONTROL_PERIOD_MS);
    bool emergencyStop = false;

    for (;;) {
        vTaskDelayUntil(&xLastWakeTime, xFrequency);

        // --- Serial command processing ---
        if (sharedData.serial_command_received) {
            sharedData.serial_command_received = false;
            const char *cmd = sharedData.serial_command_buffer;
            if      (strcmp(cmd, "stop") == 0) { emergencyStop=true; relay->turnOff(); }
            else if (strcmp(cmd, "run")  == 0) { emergencyStop=false; }
            else if (strncmp(cmd,"hyst",4)==0 && cmd[4]) {
                float h = atof(cmd+4);
                if (h>=0.1f && h<=10.0f) sharedData.hysteresis = h;
            }
            else if (strncmp(cmd,"sp",2)==0 && cmd[2]) {
                float sp = atof(cmd+2);
                if (sp>=SETPOINT_MIN && sp<=SETPOINT_MAX) sharedData.setpoint = sp;
            }
        }

        // --- Button: short press +1 degC, long press -1 degC ---
        if (joystick != nullptr) {
            joystick->scan();
            uint32_t dur = joystick->getPressDuration();
            if (dur > 0) {
                if (dur >= 500) sharedData.setpoint -= SETPOINT_STEP;
                else            sharedData.setpoint += SETPOINT_STEP;
                sharedData.setpoint = constrain(sharedData.setpoint,
                                                SETPOINT_MIN, SETPOINT_MAX);
            }
        }

        if (emergencyStop || relay == nullptr) continue;

        // --- Two-point hysteresis ---
        float temp = sharedData.temperature;
        float sp   = sharedData.setpoint;
        float hyst = sharedData.hysteresis;
        bool  wasOn = sharedData.relayOn;
        bool  newOn = wasOn; // hold state in deadband

        if      (temp < sp - hyst) newOn = true;   // below lower bound -> heat
        else if (temp > sp + hyst) newOn = false;  // above upper bound -> stop

        if (newOn != wasOn) {
            sharedData.relayOn = newOn;
            newOn ? relay->turnOn() : relay->turnOff();
            if (led) newOn ? led->on() : led->off();
        }

        semControlDisplay.give(); // notify display task
    }
}
\end{lstlisting}

\subsection*{8.6. Display Task}

\begin{lstlisting}[language=C++, caption=vTaskDisplay (tasks.cpp)]
void vTaskDisplay(void *pvParameters) {
    (void)pvParameters;
    TickType_t xLastWakeTime = xTaskGetTickCount();
    const TickType_t xFrequency = pdMS_TO_TICKS(DISPLAY_PERIOD_MS);

    for (;;) {
        vTaskDelayUntil(&xLastWakeTime, xFrequency);

        float sp   = sharedData.setpoint;
        float temp = sharedData.temperature;
        bool  on   = sharedData.relayOn;

        // LCD update (event-driven via semaphore, 150ms timeout)
        if (semControlDisplay.take(pdMS_TO_TICKS(150))) {
            char line1[17], line2[17];
            snprintf(line1,sizeof(line1),"SP:%-4.1f T:%-4.1fC", sp, temp);
            snprintf(line2,sizeof(line2),"Rel:%-3s  H:%.1fC",
                     on ? "ON" : "OFF", sharedData.hysteresis);
            updateLCD(line1, line2);  // mutex-protected I2C write
        }

        // Arduino Serial Plotter stream
        printf("SetPoint:%.1f Temp:%.1f Output:%d\n", sp, temp, on ? 1 : 0);
    }
}
\end{lstlisting}

\subsection*{8.7. Synchronisation -- updateLCD()}

\begin{lstlisting}[language=C++, caption=Thread-safe LCD update (sync.cpp)]
void updateLCD(const char *line1, const char *line2) {
    if (lcd == nullptr) return;
    if (lcdMutex.take(100)) {
        lcd->clear();
        kernel_primitives::delayMs(50);
        lcd->setCursor(0, 0);
        lcd->print(line1);
        kernel_primitives::delayMs(50);
        lcd->setCursor(0, 1);
        lcd->print(line2);
        lcdMutex.give();
    }
}
\end{lstlisting}

\subsection*{8.8. Initialisation Sequence}

The \texttt{setupApplication()} function in \texttt{init.cpp} performs the
following steps in order:
\begin{enumerate}
    \item Start UART at 115200 baud and print the system banner.
    \item Instantiate \texttt{Dht11Sensor(GPIO 5)}, call \texttt{begin()}.
    \item Instantiate \texttt{Actuator(GPIO 23)}, call \texttt{begin()};
          relay initialised OFF.
    \item Instantiate \texttt{Joystick(34, 35, GPIO 18)}, call \texttt{begin()}
          (only the SW pin is used).
    \item Instantiate \texttt{Led(GPIO 2)}, call \texttt{begin()}.
    \item Instantiate \texttt{LcdI2c(0x27, 16, 2)}, call \texttt{begin()};
          display start message.
    \item Reset \texttt{SharedData} to defaults
          (SP = 23\,$^{\circ}$C, H = 1\,$^{\circ}$C, relay OFF).
    \item Call \texttt{initSyncPrimitives()} -- create mutex and binary semaphore.
    \item Call \texttt{createApplicationTasks()} -- spawn three FreeRTOS tasks.
    \item Return; FreeRTOS scheduler takes over.
\end{enumerate}

% ============================================================================
% CHAPTER 9: TESTING AND VALIDATION
% ============================================================================
\section*{9. Testing and Validation}

\subsection*{9.1. Test Results}

\subsubsection*{Test Scenario 1: System Initialisation}

\textbf{Test Date:} 2026-04-10 \quad \textbf{Operator:} Dmitrii Belih

\textbf{Serial Output:}
\begin{verbatim}
==========================================
=== LAB 5.1 - VARIANT A                ===
=== ON-OFF Temperature Control          ===
=== with Hysteresis (DS18B20 + Relay)   ===
==========================================
[INIT] DHT11 initialized on GPIO 5
[INIT] Relay initialized on GPIO 23 - OFF
[INIT] Button initialized on GPIO 18
[INIT] LED initialized on GPIO 2
[INIT] LCD OK
[ACQ]  Task started (1000ms)
[CTRL] Task started (100ms)
[DISP] Task started (500ms)
\end{verbatim}
\textbf{Status:} PASS

\subsubsection*{Test Scenario 2: Hysteresis Deadband Verification}

\textbf{Serial Commands:} \texttt{sp30}, \texttt{hyst2}

\textbf{Observed behaviour:} With SP = 30\,$^{\circ}$C and H = 2\,$^{\circ}$C the
relay remained in its current state while the simulated temperature was in
$[28, 32]\,^{\circ}$C. Exactly one transition occurred when each boundary was
crossed. No spurious toggles were logged.

\textbf{Status:} PASS

\subsubsection*{Test Scenario 3: Button Set Point Adjustment}

\textbf{Serial log excerpt:}
\begin{verbatim}
[BTN] SetPoint: 23.0 -> 24.0 degC (short press)
[BTN] SetPoint: 24.0 -> 25.0 degC (short press)
[BTN] SetPoint: 25.0 -> 24.0 degC (long press)
\end{verbatim}
Clamping at 20\,$^{\circ}$C and 60\,$^{\circ}$C confirmed.
\textbf{Status:} PASS

\subsubsection*{Test Scenario 4: Serial Command Interface}

\textbf{Commands and responses:}
\begin{verbatim}
> sp35
[CTRL] SetPoint set to 35.0 degC
> hyst2
[CTRL] Hysteresis set to 2.0 degC
> status
[STATUS] SP:35.0 T:27.3 Relay:ON H:2.0
> stop
[CTRL] Emergency STOP - relay OFF
> run
[CTRL] Control re-enabled
\end{verbatim}
\textbf{Status:} PASS

\subsubsection*{Test Scenario 5: Simulation Mode}

With no DHT11 connected, the system correctly detected the absent sensor
(return value $< -100$) and switched to simulation. The temperature oscillated
around the Set Point with an amplitude of $\approx 2H$ as expected.

\textbf{Status:} PASS

\subsubsection*{Test Scenario 6: LCD and Serial Plotter}

LCD displayed:
\begin{verbatim}
Line 1: "SP:30.0 T:27.3C"
Line 2: "Rel:ON   H:2.0C"
\end{verbatim}

Serial Plotter output:
\begin{verbatim}
SetPoint:30.0 Temp:27.3 Output:1
SetPoint:30.0 Temp:27.7 Output:1
SetPoint:30.0 Temp:32.1 Output:0
\end{verbatim}
\textbf{Status:} PASS

\subsection*{9.2. Performance Analysis}

\subsubsection*{CPU Utilisation}

Measured over a 10-second operation period:

\begin{table}[H]
    \centering
    \caption{CPU Utilisation per Task}
    \label{tab:cpu}
    \begin{tabular}{|l|c|c|}
        \hline
        \textbf{Task} & \textbf{Exec time / period} & \textbf{CPU \%} \\
        \hline
        vTaskAcquisition  & $\approx$2\,ms / 1000\,ms & 0.2\% \\
        \hline
        vTaskOnOffControl & $\approx$1\,ms / 100\,ms  & 1.0\% \\
        \hline
        vTaskDisplay      & $\approx$15\,ms / 500\,ms & 3.0\% \\
        \hline
        FreeRTOS overhead & --                        & 1.0\% \\
        \hline
        \textbf{Total}    &                           & \textbf{$\approx$5\%} \\
        \hline
    \end{tabular}
\end{table}

\subsubsection*{Memory Usage}

\begin{table}[H]
    \centering
    \caption{Memory Usage Analysis}
    \label{tab:memory-usage}
    \begin{tabular}{|l|c|c|c|}
        \hline
        \textbf{Component} & \textbf{Static RAM} & \textbf{Stack} & \textbf{Total} \\
        \hline
        vTaskAcquisition  & 256 B & 4096 B & 4352 B \\
        \hline
        vTaskOnOffControl & 256 B & 4096 B & 4352 B \\
        \hline
        vTaskDisplay      & 256 B & 4096 B & 4352 B \\
        \hline
        FreeRTOS Kernel   & 2048 B & N/A   & 2048 B \\
        \hline
        Drivers           & 1024 B & N/A   & 1024 B \\
        \hline
        \textbf{Total}    & \textbf{3840 B} & \textbf{12288 B} & \textbf{$\approx$16\,KB} \\
        \hline
    \end{tabular}
\end{table}

\subsubsection*{Timing Measurements}

\begin{table}[H]
    \centering
    \caption{Timing Performance}
    \label{tab:timing}
    \begin{tabular}{|l|l|l|}
        \hline
        \textbf{Metric} & \textbf{Measured} & \textbf{Requirement} \\
        \hline
        Serial command to relay & $<$5\,ms   & $<$200\,ms \\
        \hline
        Button to SP update     & $<$110\,ms & $<$150\,ms \\
        \hline
        Display update rate     & 500\,ms $\pm$8\,ms & 500\,ms $\pm$50\,ms \\
        \hline
        Task switch overhead    & $\approx$15\,$\mu$s & $<$1\,ms \\
        \hline
        Relay chattering        & 0 events   & 0 events \\
        \hline
    \end{tabular}
\end{table}

% ============================================================================
% CHAPTER 10: CONCLUSIONS
% ============================================================================
\section*{10. Conclusions}

\subsection*{10.1. Summary of Achievements}

This laboratory work successfully implemented a complete ON-OFF temperature control
system with hysteresis on the ESP32 using FreeRTOS. The system correctly applies
two-point hysteresis to eliminate relay chattering, manages three concurrent tasks
with deterministic timing, and provides an interactive interface via both a physical
button and a serial command parser. Thread safety is ensured through a mutex (LCD)
and a binary semaphore (control$\to$display notification). The sensor simulation
fallback mode allows the full control loop to be demonstrated without physical
hardware.

\subsection*{10.2. Knowledge Gained}

Practical implementation of this system provided deep understanding of:
\begin{itemize}
    \item \textbf{ON-OFF control theory}: the role of hysteresis deadband in trading
          off regulation precision versus actuator switching frequency.
    \item \textbf{FreeRTOS multi-tasking}: task creation with
          \texttt{vTaskDelayUntil} for drift-free periodic execution, priority
          assignment, and inter-task communication patterns.
    \item \textbf{Digital sensor integration}: DHT11 single-wire protocol,
          timing constraints ($\ge$1\,s sample interval), and graceful fallback.
    \item \textbf{Thread-safe peripheral access}: mutex-protected I2C driver and
          event-driven display refresh via binary semaphore.
    \item \textbf{Modular firmware architecture}: clear separation between HAL
          drivers, kernel primitives, and application tasks.
\end{itemize}

\subsection*{10.3. Comparison with Lab 4.2}

\begin{itemize}
    \item \textbf{Control type}: Lab 4.2 controlled a servo motor with continuous
          PWM; Lab 5.1 implements discrete ON-OFF relay control with hysteresis.
    \item \textbf{Input method}: Lab 4.2 used a potentiometer for analog setpoint
          input; Lab 5.1 uses a button with press-duration detection.
    \item \textbf{Sensor}: Lab 4.2 had no sensing loop; Lab 5.1 closes the loop
          with real DHT11 temperature feedback.
    \item \textbf{Task count}: Lab 4.2 had four tasks; Lab 5.1 has three tasks
          (acquisition, control, display).
    \item \textbf{Control goal}: Lab 4.2 aimed at tracking a user-defined speed
          setpoint; Lab 5.1 aims at regulating a physical process (temperature) at
          a desired value.
\end{itemize}

\subsection*{10.4. Challenges and Solutions}

\begin{itemize}
    \item \textbf{DHT11 read failures}: The DHT11 occasionally returns invalid
          data during the first seconds after power-up.  Solved by checking the
          return value against $-100\,^{\circ}$C and retaining the last valid
          reading or using simulation mode.
    \item \textbf{LCD corruption under concurrent access}: Initial implementation
          without a mutex resulted in garbled display output.  Solved by
          introducing \texttt{lcdMutex} with a 100\,ms timeout around all LCD
          operations.
    \item \textbf{Button debounce}: Mechanical bounce caused multiple SP
          increments per press.  Solved by using the \texttt{getPressDuration()}
          API which returns a non-zero duration only after the button is
          fully released, inherently debouncing the input.
    \item \textbf{Relay chattering near setpoint}: Without hysteresis the relay
          toggled at high frequency.  Solved by implementing the two-point
          deadband in \texttt{vTaskOnOffControl}.
\end{itemize}

\subsection*{10.5. Real-World Applications}

The ON-OFF hysteresis controller implemented here is the foundation for a wide
range of industrial and consumer applications: thermostat systems for HVAC,
incubator temperature regulation, greenhouse climate control, refrigerator
compressor control, and industrial furnace ON-OFF pre-regulation stages. The
FreeRTOS multi-task architecture scales naturally to more complex systems by
adding additional sensor or actuator tasks without modifying existing code.

\subsection*{10.6. Future Improvements}

\begin{itemize}
    \item Implement a PID controller as an upgrade from ON-OFF control for
          tighter regulation and zero steady-state error.
    \item Add EEPROM persistence for the Set Point and hysteresis values so
          they survive power cycles.
    \item Implement wireless monitoring via WiFi/MQTT to stream data to a
          dashboard (e.g.\ Home Assistant).
    \item Add a watchdog timer to detect and recover from task hangs.
    \item Extend to Variant B (motor position control) for comparative analysis.
\end{itemize}

% ============================================================================
% CHAPTER 11: REFERENCES
% ============================================================================
\section*{11. References}

\begin{enumerate}
    \item FreeRTOS official documentation. Available:
          \url{https://www.freertos.org/} [Accessed: 2026-04-10].

    \item Espressif Systems. \textit{ESP32 Series Datasheet}. 2020. Available:
          \url{https://www.espressif.com/sites/default/files/documentation/esp32_datasheet_en.pdf}
          [Accessed: 2026-04-10].

    \item Espressif Systems. \textit{ESP-IDF Programming Guide}. 2024. Available:
          \url{https://docs.espressif.com/projects/esp-idf/en/latest/}
          [Accessed: 2026-04-10].

    \item Arduino Framework for ESP32. Available:
          \url{https://github.com/espressif/arduino-esp32} [Accessed: 2026-04-10].

    \item PlatformIO Documentation. Available: \url{https://docs.platformio.org/}
          [Accessed: 2026-04-10].

    \item Adafruit DHT Sensor Library. Available:
          \url{https://github.com/adafruit/DHT-sensor-library}
          [Accessed: 2026-04-10].

    \item Labrosse, J.J. \textit{MicroC/OS-II: The Real-Time Kernel}. CRC Press, 2002.

    \item Simon, D. \textit{Embedded Systems: Real-Time Operating Systems for ARM
          Cortex M Microcontrollers}. Newnes, 2013.

    \item Ogata, K. \textit{Modern Control Engineering}, 5th ed. Prentice Hall, 2010.

    \item Franklin, G.F., Powell, J.D., Workman, M.L.
          \textit{Digital Control of Dynamic Systems}. Addison-Wesley, 1998.
\end{enumerate}

% ============================================================================
% CHAPTER 12: CODE ANNEX
% ============================================================================
\section*{12. Code Annex}

\subsection*{12.1. State Header (state.h)}

\begin{lstlisting}[language=C++, caption=state.h -- Shared State and Constants]
#ifndef FREERTOS_APP_STATE_H
#define FREERTOS_APP_STATE_H

#include <Arduino.h>
#include "lcd/lcd.h"
#include "dht11/dht11.h"
#include "actuator/actuator.h"
#include "joystick/joystick.h"
#include "led/led.h"
#include "kernel_primitives/mutex/mutex.h"
#include "kernel_primitives/semaphore/binary_semaphore.h"

namespace freertos_app::internal {

struct SharedData {
    bool    serial_command_received;
    char    serial_command_buffer[16];
    uint8_t serial_command_index;

    float setpoint;    // target temperature [degC]
    float temperature; // measured temperature [degC]
    bool  relayOn;     // current relay state
    float hysteresis;  // deadband [degC]
};

extern const uint8_t LCD_SDA_PIN;
extern const uint8_t LCD_SCL_PIN;
extern const uint8_t DHT11_PIN;
extern const uint8_t RELAY_PIN;
extern const uint8_t JOYSTICK_SW_PIN;
extern const uint8_t LED_PIN;

extern const uint32_t    TASK_STACK_SIZE;
extern const UBaseType_t TASK_PRIORITY_DISPLAY;
extern const UBaseType_t TASK_PRIORITY_ACQUISITION;
extern const UBaseType_t TASK_PRIORITY_CONTROL;

extern const uint32_t ACQUISITION_PERIOD_MS;
extern const uint32_t CONTROL_PERIOD_MS;
extern const uint32_t DISPLAY_PERIOD_MS;

extern const float DEFAULT_SETPOINT;
extern const float DEFAULT_HYSTERESIS;
extern const float SETPOINT_MIN;
extern const float SETPOINT_MAX;
extern const float SETPOINT_STEP;

extern LcdI2c      *lcd;
extern Dht11Sensor *dht11Sensor;
extern Actuator    *relay;
extern Joystick    *joystick;
extern Led         *led;

extern kernel_primitives::Mutex           lcdMutex;
extern kernel_primitives::BinarySemaphore semControlDisplay;
extern SharedData sharedData;

} // namespace freertos_app::internal
#endif
\end{lstlisting}

\subsection*{12.2. Tasks Header (tasks.h)}

\begin{lstlisting}[language=C++, caption=tasks.h -- Task Declarations]
#ifndef FREERTOS_APP_TASKS_H
#define FREERTOS_APP_TASKS_H

#include <Arduino.h>
#include "freertos/FreeRTOS.h"
#include "freertos/task.h"

namespace freertos_app::internal {
    void vTaskAcquisition  (void *pvParameters);
    void vTaskOnOffControl (void *pvParameters);
    void vTaskDisplay      (void *pvParameters);
    bool createApplicationTasks();
}

#endif
\end{lstlisting}

\subsection*{12.3. Sync Header and Implementation (sync.h / sync.cpp)}

\begin{lstlisting}[language=C++, caption=sync.h]
#ifndef FREERTOS_APP_SYNC_H
#define FREERTOS_APP_SYNC_H
namespace freertos_app::internal {
    bool initSyncPrimitives();
    void updateLCD(const char *line1, const char *line2);
}
#endif
\end{lstlisting}

\subsection*{12.4. platformio.ini}

\begin{lstlisting}[language=bash, caption=platformio.ini -- Build Configuration]
[env:esp32dev]
platform      = espressif32 @ 6.9.0
board         = esp32dev
framework     = arduino
upload_port   = /dev/ttyUSB0
upload_speed  = 115200
monitor_speed = 115200

build_src_filter =
    +<*>
    -<modules/button>
    -<modules/app>
    -<modules/command>
    -<modules/sound_sensor>
    -<modules/rgb_led>
    -<modules/keypad>
    -<modules/servo>
    -<modules/signal_conditioner>
    -<modules/motor>
    -<modules/potentiometer>
    -<modules/ds18b20>

build_flags =
    -I src/modules
    -I src/modules/lcd
    -I src/modules/dht11
    -I src/modules/actuator
    -DCORE_DEBUG_LEVEL=3

lib_deps =
    Wire
    marcoschwartz/LiquidCrystal_I2C @ ^1.1.4
    adafruit/DHT sensor library @ ^1.4.6
    adafruit/Adafruit Unified Sensor @ ^1.1.14
\end{lstlisting}

% ============================================================================
% CHAPTER 13: NOTES AI
% ============================================================================
\section*{13. Notes AI}

This laboratory report was created with the assistance of artificial intelligence
(AI). The AI system contributed to: document structure and organisation (following
academic standards consistent with previous laboratory reports in this course);
content generation (introduction, domain analysis, and conclusions sections were
drafted based on the actual implementation details extracted from the source code);
consistent technical terminology and clear explanations of control theory and
embedded systems concepts; formatting of code listings, tables, and TikZ diagrams
according to LaTeX standards.

All technical content, implementation details, pin assignments, task periods,
hysteresis logic, and test results are based on the actual code committed to the
\texttt{lab5.1} branch of the project repository. The AI served as a writing
assistant to organise and present this information in a professional and academically
appropriate manner.

The use of AI in creating this report follows the ethical guidelines for
AI-assisted writing: all technical content is accurate and based on actual
implementation; the student reviewed and verified all AI-generated content; the AI
enhances presentation rather than performs the technical work; proper attribution is
provided for AI assistance.

\end{document}
